\documentclass[12pt]{article}
\usepackage{amsmath,amssymb,geometry,hyperref,setspace}
\geometry{margin=1in}
\usepackage[utf8]{inputenc}
\usepackage{booktabs}
\usepackage{longtable}

\title{\textbf{THE LIMIT TRIAD: A MINIMAL CONSTRAINT BASIS FOR COHERENT SYSTEMS}}
\author{}
\date{}

\begin{document}

\maketitle

\begin{abstract}
\noindent We propose that all coherent systems—formal, physical, cognitive, institutional—are constrained by exactly three irreducible limits: Gödel (internal proof), Tarski (internal truth definition), and MEL (identity under transformation). We demonstrate their independence, show that entropy and other candidate principles reduce to combinations of these three, and present a falsifiable physical instantiation (the MEL Equation) that separates the limit principle from its domain-specific applications. Stress tests across eight domains reveal no counterexamples. The result is a minimal constraint basis for any system capable of self-reference and persistence.
\end{abstract}

\section{INTRODUCTION}

\subsection{The Proliferation Problem}

Modern science and philosophy confront a peculiar embarrassment of riches: hundreds of principles, each claiming fundamental status within its domain. Physicists invoke conservation laws, thermodynamic imperatives, and quantum constraints. Computer scientists reference computational complexity bounds and the halting problem. Economists deploy equilibrium conditions and impossibility theorems. Cognitive scientists appeal to information processing limits and memory constraints. Institutional theorists cite Arrow's theorem, principal-agent problems, and constitutional design principles.

Each principle appears fundamental within its silo. Yet this proliferation raises an uncomfortable question: if all these constraints are truly fundamental, why do they fragment across domains rather than unifying? The conservation of energy matters in physics but seems irrelevant to formal logic. Gödel's incompleteness theorems constrain mathematics but appear orthogonal to biological evolution. The second law of thermodynamics governs closed physical systems but offers little guidance for institutional design.

This fragmentation suggests either that we have discovered hundreds of independent fundamental constraints—a conceptually unsatisfying position—or that many of these principles derive from a smaller set of deeper limits. The latter possibility motivates our investigation.

\subsection{The Compression Question}

We ask: what is the \textit{minimum} set of constraints that apply to \textit{all} coherent systems, regardless of domain?

This question draws inspiration from Kolmogorov's approach to complexity: seek the shortest program that generates the observed phenomenon. Applied to principles rather than data, we seek the shortest generative description of the constraint space itself. If hundreds of domain-specific principles can be derived from a handful of universal limits, we achieve not merely elegance but genuine compression of our conceptual framework.

The search for such compression must satisfy several requirements. First, the proposed limits must be demonstrably independent—no two should collapse into one. Second, they must be jointly necessary—removing any one should either permit contradictions or leave some coherent systems unconstrained. Third, they must be irreducible—no fourth limit should exist that cannot be derived from the first three. Fourth, they must apply universally across formal, physical, cognitive, and institutional domains.

\subsection{Claim}

We propose that exactly three limits satisfy these criteria:

\textbf{Gödel's Limit:} No sufficiently powerful formal system can prove all truths about itself using only internal rules. Any system capable of self-reference faces necessary incompleteness in its capacity for self-justification.

\textbf{Tarski's Limit:} No sufficiently expressive language can define its own truth predicate without contradiction. Any system capable of semantic self-reference requires external metalanguage for meaning.

\textbf{MEL's Limit:} No system can transform indefinitely without preserving traces of its origin and boundary conditions. Any system capable of temporal persistence faces necessary degradation when change outpaces provenance.

These three limits are irreducible, independent, and jointly necessary. Everything else—including entropy, conservation laws, computational complexity bounds, and institutional impossibility theorems—derives from combinations of these three fundamental constraints.

This is a strong claim. It asserts that the bewildering diversity of principles across domains reduces to three universal limits governing proof, meaning, and persistence. The remainder of this paper defends this claim through rigorous analysis.

\subsection{Paper Structure}

Section 2 establishes necessary definitions, clarifying what we mean by ``system,'' ``self-reference,'' ``coherence,'' and ``limit principle.'' Section 3 presents the three limits in canonical form. Section 4 proves their mutual independence. Section 5 demonstrates how entropy and other candidate principles reduce to these three. Section 6 stress-tests the framework across eight domains, analyzing attempted counterexamples. Section 7 resolves the apparent paradox between MEL as an unfalsifiable limit and the MEL Equation as a falsifiable physics prediction. Section 8 addresses why exactly three limits are necessary and sufficient. Section 9 explores implications for formal systems, institutions, AI alignment, physics, and philosophy of science. Section 10 identifies open questions. Section 11 concludes.

\section{DEFINITIONS}

\subsection{What Is a System?}

We define a system as any entity possessing:

\begin{enumerate}
    \item \textbf{Boundary conditions} that distinguish inside from outside
    \item \textbf{Internal representation} of states or configurations
    \item \textbf{Capacity for transformation} from one state to another
\end{enumerate}

This minimal definition deliberately spans multiple domains. A formal system (axioms, rules, theorems) qualifies. So does a physical system (boundary, microstates, dynamics). Cognitive systems (sensory boundaries, mental representations, thought processes), institutional systems (membership criteria, internal rules, policy evolution), and computational systems (input/output boundaries, memory states, program execution) all satisfy these criteria.

The definition excludes entities that lack clear boundaries (unbounded fields), internal structure (point particles under some interpretations), or transformation capacity (static eternal objects). This exclusion is intentional: our limits apply to systems capable of self-reference and persistence, which require the three features listed above.

\subsection{What Is Self-Reference?}

Self-reference occurs when a system's representation includes reference to itself. This takes three distinct forms:

\textbf{Formal self-reference} occurs when a system's rules permit statements about provability within that system. Gödel's construction shows how sufficiently powerful arithmetic can encode ``this statement is not provable in this system.'' The system refers to its own proof-generating capacity.

\textbf{Semantic self-reference} occurs when a language includes predicates that apply to sentences in that language. The liar paradox (``this sentence is false'') demonstrates semantic self-reference. The system refers to its own truth-assigning capacity.

\textbf{Temporal self-reference} occurs when a system at time $t_2$ claims identity with the system at time $t_1$ despite transformation. A person who says ``I am the same individual who existed yesterday'' engages in temporal self-reference. The system refers to its own persistence across change.

These three forms of self-reference correspond to three potential failure modes. Formal self-reference can fail through incompleteness—the system cannot prove all truths about itself. Semantic self-reference can fail through paradox—the system cannot coherently assign truth values to all self-referential statements. Temporal self-reference can fail through identity collapse—the system cannot maintain continuity across transformation.

Each failure mode corresponds to one of our three limits.

\subsection{What Is Coherence?}

A coherent system maintains:

\begin{enumerate}
    \item \textbf{Internal consistency} (no contradictions)
    \item \textbf{Meaningful reference} (truth conditions are well-defined)
    \item \textbf{Identity across transformation} (continuity through change)
\end{enumerate}

Coherence breakdown occurs when any one of these three conditions fails. A system that generates contradictions loses internal consistency. A system that cannot assign truth values to its own statements loses meaningful reference. A system that transforms beyond recognition loses identity.

Note that coherence does not require perfection. A system can be incomplete (Gödel) while remaining coherent. It can require external grounding for truth (Tarski) while remaining coherent. It can gradually drift from its origins (MEL) while remaining coherent—up to a point. The limits describe the boundaries of coherence, not barriers to functional operation.

\subsection{What Is a Limit Principle?}

A limit principle is a constraint that cannot be removed without producing one of three consequences:

\begin{enumerate}
    \item \textbf{Contradiction} (the system becomes logically inconsistent)
    \item \textbf{Collapse} (the system loses capacity for self-reference)
    \item \textbf{Trivialization} (the system becomes incapable of meaningful distinctions)
\end{enumerate}

Crucially, limit principles are definitional rather than empirical. They describe what is possible for systems of certain types, not what happens to be true in our universe. Gödel's incompleteness is not an observation about human limitations or computational resources—it is a theorem about what formal systems can and cannot do. Similarly, Tarski's undefinability is not a contingent fact about our languages—it is a necessary feature of sufficiently expressive semantic frameworks.

This definitional status distinguishes limit principles from empirical laws. The second law of thermodynamics might be violated in sufficiently unusual physical configurations; Gödel's incompleteness cannot be violated by any formal system meeting the theorem's conditions. Limit principles mark the boundaries of logical possibility, not physical actuality.

\section{THE THREE LIMITS}

\subsection{Gödel's Limit}

\textbf{Canonical Statement:} In any consistent formal system sufficiently powerful to express arithmetic, there exist true statements that cannot be proven within the system.

Kurt Gödel's 1931 incompleteness theorems demonstrated this for formal arithmetic, but the principle generalizes far beyond mathematics. The core insight concerns systems capable of encoding statements about their own proof-generating processes. Any such system faces a necessary trade-off: either it is incomplete (containing unprovable truths) or inconsistent (capable of proving falsehoods).

The standard interpretation focuses on formal provability. But the deeper constraint concerns \textit{self-justification}. A system cannot fully justify itself using only its internal resources. External reference—axioms accepted without proof, oracles consulted for undecidable questions, grounding in metalogical frameworks—

\end{document}
